\documentclass[11pt]{article}
\usepackage{lmodern}
\usepackage[T1]{fontenc}
\usepackage[letterpaper,margin=0.8in]{geometry}
\usepackage{xcolor}
\usepackage{booktabs}
\usepackage{tabularx}
\usepackage{longtable}
\usepackage{enumitem}
\usepackage{array}
\usepackage[hidelinks]{hyperref}

\setlist{itemsep=1pt,topsep=3pt}
\newcommand{\br}[1]{\texttt{#1}}
\newcommand{\h}[1]{\texttt{#1}}
\newcolumntype{L}{>{\raggedright\arraybackslash}X}

\title{PACMAN Hijinks Firmware and Software (2.X)\\[2pt]
  \large Branches, releases and bitwise reproduction}
\author{Maintained by the PACMAN agents (augie)}
\date{2026-10-31}

\begin{document}
\maketitle
\sloppy

\section{Purpose and scope}

Hijinks (2.X) is the firmware and server written in 2024 for the 10-tile PACMAN 1v5 test stands at LBNL and UC
Davis, between the legacy firmware (1.X, companion note) and zipline (3.X).  No 2.X firmware was committed in 2024:
it lived only in Vivado projects.  The record is therefore the project archives, the released bundles, and two
branches of \h{mulhearn/pacman} rebuilt from them.  This note describes the branches, the releases, and how each
released artifact is reproduced from the repository.  Bitstreams are compared by the SHA-256 of their configuration
body (first eight hex digits quoted).

\section{Branches}

\begin{itemize}
\item \br{working-on-hijinks-firmware} (da3f75c): firmware only, linear, on the SYNC FIX build step of the hw1v5
  legacy line (3155d4a), so every commit builds by construction.  Each 2.X version is one commit restored from its
  project archive, with the SYNC FIX build fixes applied to the archive's exported block design the same way.
\item \br{working-on-hijinks-software} (af73a4b): the 2024 pacman-server history as committed (server 2.1 to 2.7,
  2024-10-02 to 11-08), unchanged, linear.
\end{itemize}
Both will be superseded by a single hijinks branch of release bundles (under discussion) and then deleted.

\subsection{Firmware versions}

\begin{table}[h]
\centering\small
\begin{tabularx}{\textwidth}{@{}llllL@{}}
\toprule
Version & Commit & Date & Body & Content \\
\midrule
--- & 402413b & & & firmware only: removes the PetaLinux project and server (now on the software branch),
\h{setup.sh}, and the bare-metal apps needing an AXI GPIO the design lacks \\
2.1 & b163935 & 2024-10-05 & f87eb690 & 10 tiles, renamed I/O ports, version registers \h{0x0}--\h{0xC}, 10-bit
masks \\
2.2 & d70db3c & 2024-10-07 & 1431d8cb & hijinks constraints, 32 UARTs across 10 tiles \\
2.7 & 96d41a0 & 2024-10-16 & 6ba3a82a & UART TX inter-word pause, RX mask (2.3--2.6 were unreleased test
increments) \\
2.8 & 4b463da & 2024-10-26 & c4a85eb2 & sync/reset on the trigger line too, 32-bit RX mask, 160\,\textmu s TX gap \\
2.9 & e250cfd & 2024-10-29 & 0be6948d & sync reset pulse 15 cycles (\h{C\_SYNC\_RST\_CYCLES} 10 $\to$ 15) \\
2.10 & da3f75c & 2026-10-21 & ede22a64 & adjustable MCLK (v1.3-1.0's) on 2.8, for CERN \\
\bottomrule
\end{tabularx}
\caption{The firmware commits.  2.1--2.9 were rebuilt from the repository (Vivado 2020.2, \h{recreate\_xpr.tcl}
then \h{build\_xpr.tcl}, about 90 minutes each) and are bitwise identical to their references (2026-10-13); the
commits have since been replayed with the same trees, and each is built once more from a fresh clone of the
repository (bitwise so far: 2.1, 2.2, 2.7, 2026-11-01).  2.10 was built from this tree when released.}
\end{table}

References: 2.1 the UC Davis bundle 20241005C; 2.2 the bundles 20241007B and 20241010A (two independent bundles,
the same bitstream); 2.7 its project archive's own build (no bundle); 2.8 the LBNL bundle 20241026A; 2.9 the LBNL
bundle 20241029A; 2.10 its release 20261020A.  The project archives are \h{pacman-fw-2v1.xpr.zip} to
\h{pacman-fw-2v9.xpr.zip} (2v1, 2v2, 2v7, 2v8, 2v9).

\section{Releases}

{\small
\begin{longtable}{@{}>{\raggedright\arraybackslash}p{0.36\textwidth}>{\raggedright\arraybackslash}p{0.11\textwidth}>{\raggedright\arraybackslash}p{0.07\textwidth}>{\raggedright\arraybackslash}p{0.07\textwidth}>{\raggedright\arraybackslash}p{0.25\textwidth}@{}}
\toprule
Bundle & Kind & FW & Server & Notes \\
\midrule
\endhead
UCD 1v5\_20241005A\_ucd\_syncfix & legacy & SYNC FIX & I2C 5.2 & the branching-off point for hijinks \\
UCD 1v5\_20241005B\_ucd\_devhiijinks & development & SYNC FIX & 2.1 & server only \\
UCD 1v5\_20241005C\_ucd\_devhiijinks & development & 2.1 & 2.1 & \\
UCD 1v5\_20241007A\_ucd\_devhijinks & development & 2.1 & 2.2 & server only \\
UCD 1v5\_20241007B\_ucd\_hiijinks & released & 2.2 & 2.2 & first hijinks bundle (not successfully deployed) \\
LBNL 1v5\_20241010A\_hijinks & released & 2.2 & 2.4 & first deployed \\
LBNL 1v5\_20241026A\_hiijinks & released & 2.8 & 2.4 & \\
LBNL 1v5\_20241029A\_hijinks & unreleased & 2.9 & --- & experimental; BOOT.BIN and xsa only \\
server 2v5, 2v7 & software only & (2.8) & 2.5, 2.7 & tarballs; 2.7 is the CERN deployment \\
LBNL 1v5\_20261020A\_hijinks & released & 2.10 & 2.7 & released privately for CERN \\
\bottomrule
\end{longtable}}

The released directory names keep their ``hiijinks'' spelling.  Every bundle's \h{image.ub} is the same
(b64633c2).  The rootfs images were checked by the server binaries they carry (\h{pacman\_cmdserver},
\h{pacman\_dataserver}): each is exactly a released server tarball's.

\begin{table}[h]
\centering\small
\begin{tabular}{@{}lllll@{}}
\toprule
Server & Commit & \h{pacman\_cmdserver} & \h{pacman\_dataserver} & Shipped in \\
\midrule
2.1 & dcdbc71 & 27f457d0 & 19b33c8c & tarball 2v1 (= 05B's); rootfs of 05C \\
2.2 & 85436ea & af9df9ab & f0c56599 & tarball 2v2; rootfs of 07A, 07B \\
2.4 & 2f6f6bf & 629674af & 230cd1df & tarball 2v4; rootfs of 10A, 26A \\
2.5 & ea74a01 & 8e9c70b8 & 665e4383 & tarball 2v5 \\
2.7 & af73a4b & 61f8cb40 & 13fcbb44 & tarball 2v7 \\
\bottomrule
\end{tabular}
\caption{The released server binaries (SHA-256, first eight hex digits), each reproduced bitwise from its commit on
  \br{working-on-hijinks-software}.}\label{tab:servers}
\end{table}

\section{Reproducing the server}

{\raggedright The server is a PetaLinux 2018.3 application.  Its recipe, \h{pacman-server.bb}, is in
\h{petalinux/pkg/pacman-server}; the branch's README builds it with
\h{petalinux/scripts/build\_pacman\_server.sh} (\h{petalinux-build -c pacman-server}), and it ships as the files of
the resulting RPM.  The binaries are stripped and carry no build
dates.  Every released server was reproduced bitwise from its commit (2026-10-31, Table~\ref{tab:servers}).\par}

A plain build differs from the release in two ways, both from the build machine rather than the source:
\begin{itemize}
\item \textbf{Link order.}  The Makefile collects its sources with \h{\$(wildcard src/*.cc)}, whose order is the
  directory order of the build machine's file system.  The 2024 builds linked \h{pacman\_simulated.o},
  \h{message\_format.o}, \h{chan\_map.o} (from 2.2 on), \h{rx\_buffer.o}, \h{pacman\_i2c.o}, \h{dma.o},
  \h{tx\_buffer.o}, \h{pacman.o}, then the application.  The order was found by locating each object's code in the
  released binary, and confirmed by the 2024 build log.
\item \textbf{Build path.}  The build ID (20 bytes) and the \h{.gnu\_debuglink} checksum (4 bytes) are hashes over
  the debug information, which records the absolute path of the build's sysroot.  The 2024 builds were made in
  \h{/home/mulhearn/research/petalinux-only/pacman/petalinux/ramdisk} on vonnegut.
\end{itemize}
Each release was therefore built from its commit's recipe with the objects kept, then relinked inside the
PetaLinux 2018.3 container in the release's order, with the work directory mounted at the 2024 path, followed by
Yocto's packaging steps (\h{debugedit}, \h{objcopy -{}-only-keep-debug}, \h{-{}-add-gnu-debuglink}, \h{strip}).  The
result equals the release in every byte.  The release commits of 2.2 and 2.5, each of which spans several commits,
were pinned by the tarballs' file times (a minute or two after 85436ea and ea74a01) and confirmed by the match.
The servers in the bundles' rootfs images are the tarballs', so every released hijinks server is covered.

\section{Still open}

\begin{itemize}
\item Every firmware version built once more from a fresh clone of the repository.
\item One hijinks branch of release bundles: one commit per release, each a real merge of the firmware and
  software it brings in, confirmed by a tree check against the proven trees; annotated tags carrying each build
  result; a committed table of how each server was reproduced (commit, link order, build path, hashes).
\item Validation of 2.10 with the hardware loopback card at 5 and 10 MHz.
\end{itemize}

\end{document}
